%%%%%%%%%%%%%%%%%%%%%%%%%%%%%%%%%%%%%%%%%%%%%%%%%%%%%%%%%%%%
% 7.14 STOCHASTIC DIFFERENTIAL EQUATIONS
% The Composition of Advanced Mathematics
%%%%%%%%%%%%%%%%%%%%%%%%%%%%%%%%%%%%%%%%%%%%%%%%%%%%%%%%%%%%

\section{Stochastic Differential Equations}
\label{sec:stochastic-differential-equations}

\prerequisite{
Ordinary differential equations, stochastic processes, Brownian motion,
filtrations, martingales, It\^o integration, It\^o's formula,
conditional expectation, and numerical approximation.
}

\learningobjectives{
By the end of this section, the reader should be able to:
\begin{itemize}
    \item define a stochastic differential equation;
    \item distinguish drift and diffusion coefficients;
    \item interpret an SDE in integral form;
    \item understand strong and weak solutions;
    \item state standard existence and uniqueness conditions;
    \item solve elementary linear SDEs;
    \item derive arithmetic Brownian motion;
    \item derive geometric Brownian motion;
    \item derive the Ornstein--Uhlenbeck process;
    \item use integrating factors for linear SDEs;
    \item apply It\^o's formula to transform SDEs;
    \item compute means and variances of common SDE solutions;
    \item identify stationary distributions in basic diffusion models;
    \item understand infinitesimal generators;
    \item connect SDEs with Kolmogorov equations;
    \item recognize the Fokker--Planck equation;
    \item understand the Feynman--Kac connection;
    \item define multidimensional SDEs;
    \item understand correlated Brownian drivers;
    \item approximate SDEs using Euler--Maruyama;
    \item recognize the Milstein method;
    \item distinguish strong and weak numerical convergence;
    \item understand stability issues in numerical SDEs;
    \item recognize hitting-time and boundary-value problems;
    \item understand stochastic population and finance models;
    \item connect SDEs with diffusion processes, control theory,
          mathematical biology, physics, and finance.
\end{itemize}
}

%%%%%%%%%%%%%%%%%%%%%%%%%%%%%%%%%%%%%%%%%%%%%%%%%%%%%%%%%%%%
\subsection{From ODEs to SDEs}
\label{subsec:from-odes-to-sdes}
%%%%%%%%%%%%%%%%%%%%%%%%%%%%%%%%%%%%%%%%%%%%%%%%%%%%%%%%%%%%

An ordinary differential equation may be written

\[
\boxed{
\frac{dx}{dt}
=
b(t,x).
}
\]

Equivalently,

\[
\boxed{
dx_t
=
b(t,x_t)\,dt.
}
\]

A stochastic differential equation adds a random forcing term.

The basic It\^o form is

\[
\boxed{
dX_t
=
b(t,X_t)\,dt
+
\sigma(t,X_t)\,dB_t.
}
\]

Here:

\[
b
=
\text{drift coefficient},
\]

and

\[
\sigma
=
\text{diffusion coefficient}.
\]

%%%%%%%%%%%%%%%%%%%%%%%%%%%%%%%%%%%%%%%%%%%%%%%%%%%%%%%%%%%%
\subsection{Definition of an SDE}
\label{subsec:sde-definition}
%%%%%%%%%%%%%%%%%%%%%%%%%%%%%%%%%%%%%%%%%%%%%%%%%%%%%%%%%%%%

\begin{definitionbox}{Stochastic Differential Equation}
A stochastic differential equation driven by Brownian motion is an
equation of the form

\[
\boxed{
dX_t
=
b(t,X_t)\,dt
+
\sigma(t,X_t)\,dB_t,
}
\]

together with an initial condition

\[
\boxed{
X_0=x_0.
}
\]
\end{definitionbox}

The process \(X_t\) is unknown.

The Brownian motion \(B_t\) provides random forcing.

%%%%%%%%%%%%%%%%%%%%%%%%%%%%%%%%%%%%%%%%%%%%%%%%%%%%%%%%%%%%
\subsection{Integral Form of an SDE}
\label{subsec:sde-integral-form}
%%%%%%%%%%%%%%%%%%%%%%%%%%%%%%%%%%%%%%%%%%%%%%%%%%%%%%%%%%%%

The differential notation is shorthand.

The precise meaning is the integral equation

\[
\boxed{
X_t
=
X_0
+
\int_0^t
b(s,X_s)\,ds
+
\int_0^t
\sigma(s,X_s)\,dB_s.
}
\]

The first integral is an ordinary time integral.

The second is an It\^o stochastic integral.

%%%%%%%%%%%%%%%%%%%%%%%%%%%%%%%%%%%%%%%%%%%%%%%%%%%%%%%%%%%%
\subsection{Drift and Diffusion}
\label{subsec:sde-drift-diffusion}
%%%%%%%%%%%%%%%%%%%%%%%%%%%%%%%%%%%%%%%%%%%%%%%%%%%%%%%%%%%%

In

\[
dX_t
=
b(t,X_t)\,dt
+
\sigma(t,X_t)\,dB_t,
\]

the term

\[
b(t,X_t)\,dt
\]

describes local deterministic tendency.

The term

\[
\sigma(t,X_t)\,dB_t
\]

describes local random fluctuation.

Thus,

\[
\boxed{
\text{SDE}
=
\text{deterministic dynamics}
+
\text{random forcing}.
}
\]

%%%%%%%%%%%%%%%%%%%%%%%%%%%%%%%%%%%%%%%%%%%%%%%%%%%%%%%%%%%%
\subsection{Infinitesimal Mean and Variance}
\label{subsec:sde-infinitesimal-mean-variance}
%%%%%%%%%%%%%%%%%%%%%%%%%%%%%%%%%%%%%%%%%%%%%%%%%%%%%%%%%%%%

For a small time interval \(\Delta t\),

\[
\Delta X
\approx
b(t,X_t)\Delta t
+
\sigma(t,X_t)\Delta B.
\]

Since

\[
\mathbb{E}[\Delta B]=0,
\]

the conditional mean increment is approximately

\[
\boxed{
\mathbb{E}[\Delta X\mid X_t]
\approx
b(t,X_t)\Delta t.
}
\]

Because

\[
\operatorname{Var}(\Delta B)
=
\Delta t,
\]

the conditional variance is approximately

\[
\boxed{
\operatorname{Var}(\Delta X\mid X_t)
\approx
\sigma(t,X_t)^2\Delta t.
}
\]

%%%%%%%%%%%%%%%%%%%%%%%%%%%%%%%%%%%%%%%%%%%%%%%%%%%%%%%%%%%%
\subsection{Strong Solutions}
\label{subsec:strong-solutions-sde}
%%%%%%%%%%%%%%%%%%%%%%%%%%%%%%%%%%%%%%%%%%%%%%%%%%%%%%%%%%%%

A strong solution is constructed using a specified Brownian motion and
filtration.

Informally, \(X_t\) is built pathwise from:

\[
\boxed{
X_0
}
\]

and

\[
\boxed{
(B_s)_{0\leq s\leq t}.
}
\]

The solution must be adapted and satisfy the integral equation almost
surely.

%%%%%%%%%%%%%%%%%%%%%%%%%%%%%%%%%%%%%%%%%%%%%%%%%%%%%%%%%%%%
\subsection{Weak Solutions}
\label{subsec:weak-solutions-sde}
%%%%%%%%%%%%%%%%%%%%%%%%%%%%%%%%%%%%%%%%%%%%%%%%%%%%%%%%%%%%

A weak solution permits the probability space and Brownian motion to be
part of the construction.

The primary concern is the probability law of the resulting process
rather than a pathwise construction from a preassigned Brownian motion.

Thus:

\[
\boxed{
\text{strong solution}
=
\text{pathwise construction},
}
\]

while

\[
\boxed{
\text{weak solution}
=
\text{distributional construction}.
}
\]

%%%%%%%%%%%%%%%%%%%%%%%%%%%%%%%%%%%%%%%%%%%%%%%%%%%%%%%%%%%%
\subsection{Pathwise Uniqueness}
\label{subsec:pathwise-uniqueness}
%%%%%%%%%%%%%%%%%%%%%%%%%%%%%%%%%%%%%%%%%%%%%%%%%%%%%%%%%%%%

Suppose \(X_t\) and \(Y_t\) solve the same SDE on the same probability
space, driven by the same Brownian motion, with the same initial value.

Pathwise uniqueness means

\[
\boxed{
P
\left(
X_t=Y_t
\text{ for all }t
\right)
=
1.
}
\]

This is stronger than uniqueness only in distribution.

%%%%%%%%%%%%%%%%%%%%%%%%%%%%%%%%%%%%%%%%%%%%%%%%%%%%%%%%%%%%
\subsection{Existence and Uniqueness Conditions}
\label{subsec:sde-existence-uniqueness}
%%%%%%%%%%%%%%%%%%%%%%%%%%%%%%%%%%%%%%%%%%%%%%%%%%%%%%%%%%%%

A standard theorem assumes Lipschitz and linear-growth conditions.

Suppose there exists \(K>0\) such that

\[
\boxed{
|b(t,x)-b(t,y)|
+
|\sigma(t,x)-\sigma(t,y)|
\leq
K|x-y|
}
\]

for all relevant \(t,x,y\).

Also suppose

\[
\boxed{
|b(t,x)|^2
+
|\sigma(t,x)|^2
\leq
K^2(1+|x|^2).
}
\]

Under standard hypotheses, the SDE has a unique strong solution.

%%%%%%%%%%%%%%%%%%%%%%%%%%%%%%%%%%%%%%%%%%%%%%%%%%%%%%%%%%%%
\subsection{Why Lipschitz Conditions Matter}
\label{subsec:why-lipschitz-sde}
%%%%%%%%%%%%%%%%%%%%%%%%%%%%%%%%%%%%%%%%%%%%%%%%%%%%%%%%%%%%

The Lipschitz condition limits how rapidly the coefficients can separate
nearby states.

This prevents two solutions starting from the same point from diverging
too rapidly.

It is the stochastic analogue of classical uniqueness conditions for
ordinary differential equations.

%%%%%%%%%%%%%%%%%%%%%%%%%%%%%%%%%%%%%%%%%%%%%%%%%%%%%%%%%%%%
\subsection{Arithmetic Brownian Motion}
\label{subsec:arithmetic-brownian-motion}
%%%%%%%%%%%%%%%%%%%%%%%%%%%%%%%%%%%%%%%%%%%%%%%%%%%%%%%%%%%%

Consider

\[
\boxed{
dX_t
=
\mu\,dt
+
\sigma\,dB_t,
}
\]

with constants \(\mu\) and \(\sigma\).

Integrating directly,

\[
X_t
=
X_0
+
\mu t
+
\sigma B_t.
\]

Therefore,

\[
\boxed{
X_t
=
X_0+\mu t+\sigma B_t.
}
\]

This is arithmetic Brownian motion.

%%%%%%%%%%%%%%%%%%%%%%%%%%%%%%%%%%%%%%%%%%%%%%%%%%%%%%%%%%%%
\subsection{Distribution of Arithmetic Brownian Motion}
\label{subsec:arithmetic-brownian-distribution}
%%%%%%%%%%%%%%%%%%%%%%%%%%%%%%%%%%%%%%%%%%%%%%%%%%%%%%%%%%%%

Since

\[
B_t\sim N(0,t),
\]

we obtain

\[
\boxed{
X_t
\sim
N
\left(
X_0+\mu t,
\sigma^2t
\right).
}
\]

Thus,

\[
\boxed{
\mathbb{E}[X_t]
=
X_0+\mu t,
}
\]

and

\[
\boxed{
\operatorname{Var}(X_t)
=
\sigma^2t.
}
\]

%%%%%%%%%%%%%%%%%%%%%%%%%%%%%%%%%%%%%%%%%%%%%%%%%%%%%%%%%%%%
\subsection{Worked Example: Arithmetic Brownian Motion}
\label{subsec:worked-arithmetic-brownian}
%%%%%%%%%%%%%%%%%%%%%%%%%%%%%%%%%%%%%%%%%%%%%%%%%%%%%%%%%%%%

\begin{workedexamplebox}{Constant Drift and Diffusion}

Suppose

\[
dX_t
=
3\,dt
+
2\,dB_t,
\qquad
X_0=5.
\]

Integrating gives

\[
X_t
=
5+3t+2B_t.
\]

Therefore,

\begin{answerbox}
\[
\boxed{
X_t
\sim
N(5+3t,4t).
}
\]
\end{answerbox}

\end{workedexamplebox}

%%%%%%%%%%%%%%%%%%%%%%%%%%%%%%%%%%%%%%%%%%%%%%%%%%%%%%%%%%%%
\subsection{Geometric Brownian Motion}
\label{subsec:geometric-brownian-motion}
%%%%%%%%%%%%%%%%%%%%%%%%%%%%%%%%%%%%%%%%%%%%%%%%%%%%%%%%%%%%

Consider the multiplicative SDE

\[
\boxed{
dS_t
=
\mu S_t\,dt
+
\sigma S_t\,dB_t,
}
\]

with

\[
S_0>0.
\]

This is geometric Brownian motion.

It is one of the most important explicitly solvable SDEs.

%%%%%%%%%%%%%%%%%%%%%%%%%%%%%%%%%%%%%%%%%%%%%%%%%%%%%%%%%%%%
\subsection{Solving Geometric Brownian Motion}
\label{subsec:solve-geometric-brownian}
%%%%%%%%%%%%%%%%%%%%%%%%%%%%%%%%%%%%%%%%%%%%%%%%%%%%%%%%%%%%

Apply It\^o's formula to

\[
f(S_t)
=
\ln S_t.
\]

From the previous section,

\[
d(\ln S_t)
=
\left(
\mu-\frac12\sigma^2
\right)dt
+
\sigma\,dB_t.
\]

Integrating,

\[
\ln S_t
=
\ln S_0
+
\left(
\mu-\frac12\sigma^2
\right)t
+
\sigma B_t.
\]

Exponentiating,

\[
\boxed{
S_t
=
S_0
\exp
\left[
\left(
\mu-\frac12\sigma^2
\right)t
+
\sigma B_t
\right].
}
\]

%%%%%%%%%%%%%%%%%%%%%%%%%%%%%%%%%%%%%%%%%%%%%%%%%%%%%%%%%%%%
\subsection{Distribution of Geometric Brownian Motion}
\label{subsec:gbm-distribution}
%%%%%%%%%%%%%%%%%%%%%%%%%%%%%%%%%%%%%%%%%%%%%%%%%%%%%%%%%%%%

Because

\[
\ln S_t
\]

is normal,

\[
S_t
\]

is lognormal.

Specifically,

\[
\boxed{
\ln S_t
\sim
N
\left(
\ln S_0
+
\left(
\mu-\frac12\sigma^2
\right)t,
\sigma^2t
\right).
}
\]

Thus,

\[
\boxed{
S_t>0
}
\]

almost surely whenever \(S_0>0\).

%%%%%%%%%%%%%%%%%%%%%%%%%%%%%%%%%%%%%%%%%%%%%%%%%%%%%%%%%%%%
\subsection{Mean of Geometric Brownian Motion}
\label{subsec:gbm-mean}
%%%%%%%%%%%%%%%%%%%%%%%%%%%%%%%%%%%%%%%%%%%%%%%%%%%%%%%%%%%%

Using the lognormal moment formula,

\[
\mathbb{E}[S_t]
=
S_0
e^{(\mu-\sigma^2/2)t}
e^{\sigma^2t/2}.
\]

Therefore,

\[
\boxed{
\mathbb{E}[S_t]
=
S_0e^{\mu t}.
}
\]

The It\^o correction disappears from the final mean.

%%%%%%%%%%%%%%%%%%%%%%%%%%%%%%%%%%%%%%%%%%%%%%%%%%%%%%%%%%%%
\subsection{Variance of Geometric Brownian Motion}
\label{subsec:gbm-variance}
%%%%%%%%%%%%%%%%%%%%%%%%%%%%%%%%%%%%%%%%%%%%%%%%%%%%%%%%%%%%

The variance is

\[
\boxed{
\operatorname{Var}(S_t)
=
S_0^2
e^{2\mu t}
\left(
e^{\sigma^2t}-1
\right).
}
\]

Thus uncertainty increases with both time and volatility.

%%%%%%%%%%%%%%%%%%%%%%%%%%%%%%%%%%%%%%%%%%%%%%%%%%%%%%%%%%%%
\subsection{Worked Example: Geometric Brownian Motion}
\label{subsec:worked-gbm}
%%%%%%%%%%%%%%%%%%%%%%%%%%%%%%%%%%%%%%%%%%%%%%%%%%%%%%%%%%%%

\begin{workedexamplebox}{Explicit GBM Solution}

Suppose

\[
dS_t
=
0.08S_t\,dt
+
0.20S_t\,dB_t,
\qquad
S_0=100.
\]

Then

\[
S_t
=
100
\exp
\left[
\left(
0.08-\frac12(0.20)^2
\right)t
+
0.20B_t
\right].
\]

Since

\[
\frac12(0.20)^2=0.02,
\]

we obtain

\begin{answerbox}
\[
\boxed{
S_t
=
100
\exp
\left(
0.06t+0.20B_t
\right).
}
\]
\end{answerbox}

\end{workedexamplebox}

%%%%%%%%%%%%%%%%%%%%%%%%%%%%%%%%%%%%%%%%%%%%%%%%%%%%%%%%%%%%
\subsection{Linear SDEs}
\label{subsec:linear-sdes}
%%%%%%%%%%%%%%%%%%%%%%%%%%%%%%%%%%%%%%%%%%%%%%%%%%%%%%%%%%%%

A general scalar linear SDE may be written

\[
\boxed{
dX_t
=
(a_tX_t+c_t)\,dt
+
(b_tX_t+d_t)\,dB_t.
}
\]

Some linear SDEs can be solved using stochastic integrating factors.

A particularly important additive-noise example is

\[
\boxed{
dX_t
=
(aX_t+c)\,dt
+
\sigma\,dB_t.
}
\]

%%%%%%%%%%%%%%%%%%%%%%%%%%%%%%%%%%%%%%%%%%%%%%%%%%%%%%%%%%%%
\subsection{Ornstein--Uhlenbeck Process}
\label{subsec:ornstein-uhlenbeck-process}
%%%%%%%%%%%%%%%%%%%%%%%%%%%%%%%%%%%%%%%%%%%%%%%%%%%%%%%%%%%%

The Ornstein--Uhlenbeck process satisfies

\[
\boxed{
dX_t
=
\theta(\mu-X_t)\,dt
+
\sigma\,dB_t,
}
\]

where

\[
\theta>0.
\]

The drift pulls \(X_t\) toward the level \(\mu\).

Thus the process is mean-reverting.

%%%%%%%%%%%%%%%%%%%%%%%%%%%%%%%%%%%%%%%%%%%%%%%%%%%%%%%%%%%%
\subsection{Interpretation of Mean Reversion}
\label{subsec:mean-reversion-ou}
%%%%%%%%%%%%%%%%%%%%%%%%%%%%%%%%%%%%%%%%%%%%%%%%%%%%%%%%%%%%

If

\[
X_t>\mu,
\]

then

\[
\theta(\mu-X_t)<0,
\]

so the drift points downward.

If

\[
X_t<\mu,
\]

then

\[
\theta(\mu-X_t)>0,
\]

so the drift points upward.

Therefore,

\[
\boxed{
\mu
=
\text{long-run equilibrium level},
}
\]

while

\[
\boxed{
\theta
=
\text{mean-reversion speed}.
}
\]

%%%%%%%%%%%%%%%%%%%%%%%%%%%%%%%%%%%%%%%%%%%%%%%%%%%%%%%%%%%%
\subsection{Solving the Ornstein--Uhlenbeck SDE}
\label{subsec:solve-ou}
%%%%%%%%%%%%%%%%%%%%%%%%%%%%%%%%%%%%%%%%%%%%%%%%%%%%%%%%%%%%

Rewrite

\[
dX_t
+
\theta X_t\,dt
=
\theta\mu\,dt
+
\sigma\,dB_t.
\]

Use the integrating factor

\[
e^{\theta t}.
\]

Then

\[
d
\left(
e^{\theta t}X_t
\right)
=
\theta\mu e^{\theta t}\,dt
+
\sigma e^{\theta t}\,dB_t.
\]

Integrating and simplifying,

\[
\boxed{
X_t
=
\mu
+
(X_0-\mu)e^{-\theta t}
+
\sigma
\int_0^t
e^{-\theta(t-s)}
\,dB_s.
}
\]

%%%%%%%%%%%%%%%%%%%%%%%%%%%%%%%%%%%%%%%%%%%%%%%%%%%%%%%%%%%%
\subsection{Mean of the Ornstein--Uhlenbeck Process}
\label{subsec:ou-mean}
%%%%%%%%%%%%%%%%%%%%%%%%%%%%%%%%%%%%%%%%%%%%%%%%%%%%%%%%%%%%

The stochastic integral has mean zero.

Therefore,

\[
\boxed{
\mathbb{E}[X_t]
=
\mu
+
(X_0-\mu)e^{-\theta t}.
}
\]

As

\[
t\to\infty,
\]

\[
\boxed{
\mathbb{E}[X_t]
\to
\mu.
}
\]

%%%%%%%%%%%%%%%%%%%%%%%%%%%%%%%%%%%%%%%%%%%%%%%%%%%%%%%%%%%%
\subsection{Variance of the Ornstein--Uhlenbeck Process}
\label{subsec:ou-variance}
%%%%%%%%%%%%%%%%%%%%%%%%%%%%%%%%%%%%%%%%%%%%%%%%%%%%%%%%%%%%

By the It\^o isometry,

\[
\operatorname{Var}(X_t)
=
\sigma^2
\int_0^t
e^{-2\theta(t-s)}
\,ds.
\]

Therefore,

\[
\boxed{
\operatorname{Var}(X_t)
=
\frac{
\sigma^2
}{
2\theta
}
\left(
1-e^{-2\theta t}
\right).
}
\]

As \(t\to\infty\),

\[
\boxed{
\operatorname{Var}(X_t)
\to
\frac{\sigma^2}{2\theta}.
}
\]

%%%%%%%%%%%%%%%%%%%%%%%%%%%%%%%%%%%%%%%%%%%%%%%%%%%%%%%%%%%%
\subsection{Stationary Ornstein--Uhlenbeck Distribution}
\label{subsec:ou-stationary-distribution}
%%%%%%%%%%%%%%%%%%%%%%%%%%%%%%%%%%%%%%%%%%%%%%%%%%%%%%%%%%%%

The limiting distribution is Gaussian:

\[
\boxed{
X_\infty
\sim
N
\left(
\mu,
\frac{\sigma^2}{2\theta}
\right).
}
\]

Thus the Ornstein--Uhlenbeck process is one of the basic continuous-time
stationary Gaussian models.

%%%%%%%%%%%%%%%%%%%%%%%%%%%%%%%%%%%%%%%%%%%%%%%%%%%%%%%%%%%%
\subsection{Worked Example: Ornstein--Uhlenbeck Process}
\label{subsec:worked-ou}
%%%%%%%%%%%%%%%%%%%%%%%%%%%%%%%%%%%%%%%%%%%%%%%%%%%%%%%%%%%%

\begin{workedexamplebox}{Long-Run Distribution}

Suppose

\[
dX_t
=
2(10-X_t)\,dt
+
3\,dB_t.
\]

Then

\[
\theta=2,
\qquad
\mu=10,
\qquad
\sigma=3.
\]

The stationary variance is

\[
\frac{9}{4}.
\]

Therefore,

\begin{answerbox}
\[
\boxed{
X_\infty
\sim
N
\left(
10,
\frac94
\right).
}
\]
\end{answerbox}

\end{workedexamplebox}

%%%%%%%%%%%%%%%%%%%%%%%%%%%%%%%%%%%%%%%%%%%%%%%%%%%%%%%%%%%%
\subsection{General Additive Linear SDE}
\label{subsec:general-additive-linear-sde}
%%%%%%%%%%%%%%%%%%%%%%%%%%%%%%%%%%%%%%%%%%%%%%%%%%%%%%%%%%%%

Consider

\[
\boxed{
dX_t
=
(a(t)X_t+c(t))\,dt
+
\sigma(t)\,dB_t.
}
\]

An integrating factor is

\[
\boxed{
I(t)
=
\exp
\left(
-\int_0^t
a(s)\,ds
\right).
}
\]

This transforms the SDE into an equation whose drift no longer contains
the unknown process multiplicatively.

%%%%%%%%%%%%%%%%%%%%%%%%%%%%%%%%%%%%%%%%%%%%%%%%%%%%%%%%%%%%
\subsection{Transforming SDEs with It\^o's Formula}
\label{subsec:transform-sdes-ito}
%%%%%%%%%%%%%%%%%%%%%%%%%%%%%%%%%%%%%%%%%%%%%%%%%%%%%%%%%%%%

Suppose

\[
dX_t
=
b(X_t)\,dt
+
\sigma(X_t)\,dB_t
\]

and define

\[
Y_t=g(X_t).
\]

Then It\^o's formula gives

\[
\boxed{
dY_t
=
\left[
b(X_t)g'(X_t)
+
\frac12
\sigma(X_t)^2g''(X_t)
\right]dt
+
\sigma(X_t)g'(X_t)\,dB_t.
}
\]

This is the stochastic analogue of change of variables in an ODE.

%%%%%%%%%%%%%%%%%%%%%%%%%%%%%%%%%%%%%%%%%%%%%%%%%%%%%%%%%%%%
\subsection{Worked Example: Exponential Transformation}
\label{subsec:worked-sde-exponential-transform}
%%%%%%%%%%%%%%%%%%%%%%%%%%%%%%%%%%%%%%%%%%%%%%%%%%%%%%%%%%%%

\begin{workedexamplebox}{Transforming Arithmetic Brownian Motion}

Suppose

\[
dX_t
=
\mu\,dt
+
\sigma\,dB_t
\]

and define

\[
Y_t=e^{X_t}.
\]

Then

\[
g'(x)=e^x,
\qquad
g''(x)=e^x.
\]

Therefore,

\[
dY_t
=
Y_t
\left(
\mu+\frac12\sigma^2
\right)dt
+
\sigma Y_t\,dB_t.
\]

\begin{answerbox}
\[
\boxed{
dY_t
=
\left(
\mu+\frac12\sigma^2
\right)Y_t\,dt
+
\sigma Y_t\,dB_t.
}
\]
\end{answerbox}

\end{workedexamplebox}

%%%%%%%%%%%%%%%%%%%%%%%%%%%%%%%%%%%%%%%%%%%%%%%%%%%%%%%%%%%%
\subsection{Diffusion Processes}
\label{subsec:diffusion-processes}
%%%%%%%%%%%%%%%%%%%%%%%%%%%%%%%%%%%%%%%%%%%%%%%%%%%%%%%%%%%%

A diffusion is commonly modeled by an SDE of the form

\[
\boxed{
dX_t
=
b(t,X_t)\,dt
+
\sigma(t,X_t)\,dB_t.
}
\]

Under standard assumptions, such processes have continuous sample paths.

They represent systems with continuous random motion.

%%%%%%%%%%%%%%%%%%%%%%%%%%%%%%%%%%%%%%%%%%%%%%%%%%%%%%%%%%%%
\subsection{Infinitesimal Generator}
\label{subsec:sde-generator}
%%%%%%%%%%%%%%%%%%%%%%%%%%%%%%%%%%%%%%%%%%%%%%%%%%%%%%%%%%%%

For a time-homogeneous diffusion

\[
dX_t
=
b(X_t)\,dt
+
\sigma(X_t)\,dB_t,
\]

the infinitesimal generator is

\[
\boxed{
\mathcal{L}f(x)
=
b(x)f'(x)
+
\frac12
\sigma(x)^2f''(x).
}
\]

This operator describes the local evolution of expected functions of
the process.

%%%%%%%%%%%%%%%%%%%%%%%%%%%%%%%%%%%%%%%%%%%%%%%%%%%%%%%%%%%%
\subsection{Generator from It\^o's Formula}
\label{subsec:generator-from-ito}
%%%%%%%%%%%%%%%%%%%%%%%%%%%%%%%%%%%%%%%%%%%%%%%%%%%%%%%%%%%%

It\^o's formula gives

\[
df(X_t)
=
\mathcal{L}f(X_t)\,dt
+
\sigma(X_t)f'(X_t)\,dB_t.
\]

Taking conditional expectations over a small time interval yields

\[
\boxed{
\mathbb{E}
[
f(X_{t+\Delta t})-f(X_t)
\mid
X_t=x
]
\approx
\mathcal{L}f(x)\Delta t.
}
\]

This explains the term ``generator.''

%%%%%%%%%%%%%%%%%%%%%%%%%%%%%%%%%%%%%%%%%%%%%%%%%%%%%%%%%%%%
\subsection{Kolmogorov Backward Equation}
\label{subsec:kolmogorov-backward-sde}
%%%%%%%%%%%%%%%%%%%%%%%%%%%%%%%%%%%%%%%%%%%%%%%%%%%%%%%%%%%%

Let

\[
u(t,x)
=
\mathbb{E}_x
[
g(X_t)
].
\]

Under appropriate smoothness assumptions,

\[
u
\]

satisfies the backward equation

\[
\boxed{
\frac{\partial u}{\partial t}
=
\mathcal{L}u.
}
\]

For the diffusion generator,

\[
\boxed{
u_t
=
b(x)u_x
+
\frac12
\sigma(x)^2u_{xx}.
}
\]

%%%%%%%%%%%%%%%%%%%%%%%%%%%%%%%%%%%%%%%%%%%%%%%%%%%%%%%%%%%%
\subsection{Fokker--Planck Equation}
\label{subsec:fokker-planck}
%%%%%%%%%%%%%%%%%%%%%%%%%%%%%%%%%%%%%%%%%%%%%%%%%%%%%%%%%%%%

The probability density of a diffusion evolves according to a forward
equation.

Suppose

\[
p(t,x)
\]

is the density of \(X_t\).

For

\[
dX_t
=
b(X_t)\,dt
+
\sigma(X_t)\,dB_t,
\]

the Fokker--Planck equation is formally

\[
\boxed{
\frac{\partial p}{\partial t}
=
-
\frac{\partial}{\partial x}
\left(
b(x)p
\right)
+
\frac12
\frac{\partial^2}{\partial x^2}
\left(
\sigma(x)^2p
\right).
}
\]

It is also called the Kolmogorov forward equation.

%%%%%%%%%%%%%%%%%%%%%%%%%%%%%%%%%%%%%%%%%%%%%%%%%%%%%%%%%%%%
\subsection{Pure Brownian Fokker--Planck Equation}
\label{subsec:brownian-fokker-planck}
%%%%%%%%%%%%%%%%%%%%%%%%%%%%%%%%%%%%%%%%%%%%%%%%%%%%%%%%%%%%

For

\[
dX_t=dB_t,
\]

we have

\[
b=0,
\qquad
\sigma=1.
\]

Therefore,

\[
\boxed{
p_t
=
\frac12p_{xx}.
}
\]

This is the heat equation.

Thus Brownian probability densities evolve according to diffusion.

%%%%%%%%%%%%%%%%%%%%%%%%%%%%%%%%%%%%%%%%%%%%%%%%%%%%%%%%%%%%
\subsection{Stationary Densities}
\label{subsec:stationary-densities-sde}
%%%%%%%%%%%%%%%%%%%%%%%%%%%%%%%%%%%%%%%%%%%%%%%%%%%%%%%%%%%%

A stationary density

\[
p_\infty(x)
\]

does not change with time.

Therefore it solves the stationary Fokker--Planck equation

\[
\boxed{
0
=
-
\frac{d}{dx}
(bp_\infty)
+
\frac12
\frac{d^2}{dx^2}
(\sigma^2p_\infty).
}
\]

Solving this equation can reveal long-run distributions of diffusions.

%%%%%%%%%%%%%%%%%%%%%%%%%%%%%%%%%%%%%%%%%%%%%%%%%%%%%%%%%%%%
\subsection{Feynman--Kac Formula}
\label{subsec:feynman-kac-sde}
%%%%%%%%%%%%%%%%%%%%%%%%%%%%%%%%%%%%%%%%%%%%%%%%%%%%%%%%%%%%

A basic Feynman--Kac principle connects diffusion expectations with PDEs.

Suppose

\[
u(t,x)
=
\mathbb{E}_{t,x}
\left[
g(X_T)
\right].
\]

Under suitable conditions,

\[
u
\]

solves a backward PDE of the form

\[
\boxed{
u_t
+
\mathcal{L}u
=
0
}
\]

with terminal condition

\[
\boxed{
u(T,x)=g(x).
}
\]

More general versions include discounting and source terms.

%%%%%%%%%%%%%%%%%%%%%%%%%%%%%%%%%%%%%%%%%%%%%%%%%%%%%%%%%%%%
\subsection{Hitting Probabilities}
\label{subsec:sde-hitting-probabilities}
%%%%%%%%%%%%%%%%%%%%%%%%%%%%%%%%%%%%%%%%%%%%%%%%%%%%%%%%%%%%

Let

\[
\tau_A
=
\inf
\{
t\geq0:X_t\in A
\}.
\]

One may seek probabilities such as

\[
\boxed{
P_x
(
\tau_b<\tau_a
).
}
\]

For one-dimensional diffusions, such probabilities can often be found
from boundary-value problems involving the generator.

%%%%%%%%%%%%%%%%%%%%%%%%%%%%%%%%%%%%%%%%%%%%%%%%%%%%%%%%%%%%
\subsection{Generator Equation for Hitting Probabilities}
\label{subsec:generator-hitting-probabilities}
%%%%%%%%%%%%%%%%%%%%%%%%%%%%%%%%%%%%%%%%%%%%%%%%%%%%%%%%%%%%

Let

\[
h(x)
=
P_x
(
\tau_b<\tau_a
)
\]

for \(a<x<b\).

Under standard assumptions,

\[
\boxed{
\mathcal{L}h(x)=0
}
\]

with boundary conditions

\[
\boxed{
h(a)=0,
\qquad
h(b)=1.
}
\]

This converts a stochastic hitting problem into a differential-equation
problem.

%%%%%%%%%%%%%%%%%%%%%%%%%%%%%%%%%%%%%%%%%%%%%%%%%%%%%%%%%%%%
\subsection{Expected Hitting Times}
\label{subsec:expected-hitting-times-sde}
%%%%%%%%%%%%%%%%%%%%%%%%%%%%%%%%%%%%%%%%%%%%%%%%%%%%%%%%%%%%

Let

\[
u(x)
=
\mathbb{E}_x[\tau]
\]

for an exit time \(\tau\).

A common boundary-value problem is

\[
\boxed{
\mathcal{L}u=-1
}
\]

inside the domain, with

\[
\boxed{
u=0
}
\]

on the exit boundary.

This is another important stochastic/PDE connection.

%%%%%%%%%%%%%%%%%%%%%%%%%%%%%%%%%%%%%%%%%%%%%%%%%%%%%%%%%%%%
\subsection{Multidimensional SDEs}
\label{subsec:multidimensional-sdes}
%%%%%%%%%%%%%%%%%%%%%%%%%%%%%%%%%%%%%%%%%%%%%%%%%%%%%%%%%%%%

Let

\[
\mathbf{X}_t
\in
\mathbb{R}^d
\]

and let

\[
\mathbf{B}_t
\in
\mathbb{R}^m
\]

be multidimensional Brownian motion.

A general SDE is

\[
\boxed{
d\mathbf{X}_t
=
\mathbf{b}(t,\mathbf{X}_t)\,dt
+
\Sigma(t,\mathbf{X}_t)\,d\mathbf{B}_t.
}
\]

Here

\[
\mathbf{b}
\]

is a drift vector and

\[
\Sigma
\]

is a diffusion matrix.

%%%%%%%%%%%%%%%%%%%%%%%%%%%%%%%%%%%%%%%%%%%%%%%%%%%%%%%%%%%%
\subsection{Covariance Matrix of Multidimensional Noise}
\label{subsec:multidimensional-sde-covariance}
%%%%%%%%%%%%%%%%%%%%%%%%%%%%%%%%%%%%%%%%%%%%%%%%%%%%%%%%%%%%

For a small interval,

\[
d\mathbf{X}_t
\]

has instantaneous covariance matrix approximately

\[
\boxed{
\Sigma\Sigma^T\,dt.
}
\]

Thus the matrix

\[
\boxed{
A
=
\Sigma\Sigma^T
}
\]

is the diffusion covariance matrix.

%%%%%%%%%%%%%%%%%%%%%%%%%%%%%%%%%%%%%%%%%%%%%%%%%%%%%%%%%%%%
\subsection{Multidimensional Generator}
\label{subsec:multidimensional-generator-sde}
%%%%%%%%%%%%%%%%%%%%%%%%%%%%%%%%%%%%%%%%%%%%%%%%%%%%%%%%%%%%

For a smooth function

\[
f:\mathbb{R}^d\to\mathbb{R},
\]

the generator is

\[
\boxed{
\mathcal{L}f
=
\sum_{i=1}^{d}
b_i
\frac{\partial f}{\partial x_i}
+
\frac12
\sum_{i,j=1}^{d}
A_{ij}
\frac{\partial^2 f}
{\partial x_i\partial x_j}.
}
\]

Here

\[
A
=
\Sigma\Sigma^T.
\]

%%%%%%%%%%%%%%%%%%%%%%%%%%%%%%%%%%%%%%%%%%%%%%%%%%%%%%%%%%%%
\subsection{Correlated Brownian Drivers}
\label{subsec:correlated-brownian-drivers}
%%%%%%%%%%%%%%%%%%%%%%%%%%%%%%%%%%%%%%%%%%%%%%%%%%%%%%%%%%%%

Suppose

\[
dB_t^{(1)}dB_t^{(2)}
=
\rho\,dt.
\]

Then the two sources of randomness are correlated.

Such models are common when multiple stochastic quantities are driven by
shared economic, physical, or environmental effects.

%%%%%%%%%%%%%%%%%%%%%%%%%%%%%%%%%%%%%%%%%%%%%%%%%%%%%%%%%%%%
\subsection{Stochastic Logistic Growth}
\label{subsec:stochastic-logistic-growth}
%%%%%%%%%%%%%%%%%%%%%%%%%%%%%%%%%%%%%%%%%%%%%%%%%%%%%%%%%%%%

The deterministic logistic equation is

\[
\frac{dx}{dt}
=
rx
\left(
1-\frac{x}{K}
\right).
\]

A stochastic version is

\[
\boxed{
dX_t
=
rX_t
\left(
1-\frac{X_t}{K}
\right)dt
+
\sigma X_t\,dB_t.
}
\]

Here:

\[
r
=
\text{growth rate},
\]

\[
K
=
\text{carrying capacity},
\]

and

\[
\sigma
=
\text{environmental-noise intensity}.
\]

%%%%%%%%%%%%%%%%%%%%%%%%%%%%%%%%%%%%%%%%%%%%%%%%%%%%%%%%%%%%
\subsection{Mean-Reverting Finance Models}
\label{subsec:mean-reverting-finance-sdes}
%%%%%%%%%%%%%%%%%%%%%%%%%%%%%%%%%%%%%%%%%%%%%%%%%%%%%%%%%%%%

Short-term interest rates and other economic quantities are often
modeled as mean-reverting diffusions.

The Vasicek model is an Ornstein--Uhlenbeck-type SDE:

\[
\boxed{
dr_t
=
\kappa(\theta-r_t)\,dt
+
\sigma\,dB_t.
}
\]

The parameter \(\theta\) is the long-run level and \(\kappa\) controls
the reversion speed.

%%%%%%%%%%%%%%%%%%%%%%%%%%%%%%%%%%%%%%%%%%%%%%%%%%%%%%%%%%%%
\subsection{Cox--Ingersoll--Ross Model Preview}
\label{subsec:cir-preview}
%%%%%%%%%%%%%%%%%%%%%%%%%%%%%%%%%%%%%%%%%%%%%%%%%%%%%%%%%%%%

A positive mean-reverting model is

\[
\boxed{
dX_t
=
\kappa(\theta-X_t)\,dt
+
\sigma\sqrt{X_t}\,dB_t.
}
\]

This is the Cox--Ingersoll--Ross, or CIR, process.

The state-dependent diffusion

\[
\sqrt{X_t}
\]

reduces volatility near zero.

Under suitable parameter conditions, positivity is preserved.

%%%%%%%%%%%%%%%%%%%%%%%%%%%%%%%%%%%%%%%%%%%%%%%%%%%%%%%%%%%%
\subsection{Langevin Equations}
\label{subsec:langevin-equations}
%%%%%%%%%%%%%%%%%%%%%%%%%%%%%%%%%%%%%%%%%%%%%%%%%%%%%%%%%%%%

In physics, stochastic forcing is often represented through a Langevin
equation.

A simple model is

\[
\boxed{
dV_t
=
-\gamma V_t\,dt
+
\sigma\,dB_t.
}
\]

This is an Ornstein--Uhlenbeck equation for velocity.

The term

\[
-\gamma V_t
\]

represents friction, while Brownian forcing represents thermal noise.

%%%%%%%%%%%%%%%%%%%%%%%%%%%%%%%%%%%%%%%%%%%%%%%%%%%%%%%%%%%%
\subsection{Euler--Maruyama Method}
\label{subsec:euler-maruyama}
%%%%%%%%%%%%%%%%%%%%%%%%%%%%%%%%%%%%%%%%%%%%%%%%%%%%%%%%%%%%

Many SDEs do not have closed-form solutions.

The simplest numerical method is Euler--Maruyama.

For

\[
dX_t
=
b(t,X_t)\,dt
+
\sigma(t,X_t)\,dB_t,
\]

use a grid

\[
t_n=n\Delta t.
\]

Then

\[
\boxed{
X_{n+1}
=
X_n
+
b(t_n,X_n)\Delta t
+
\sigma(t_n,X_n)\Delta B_n,
}
\]

where

\[
\boxed{
\Delta B_n
\sim
N(0,\Delta t).
}
\]

%%%%%%%%%%%%%%%%%%%%%%%%%%%%%%%%%%%%%%%%%%%%%%%%%%%%%%%%%%%%
\subsection{Simulating Brownian Increments}
\label{subsec:simulate-brownian-increments}
%%%%%%%%%%%%%%%%%%%%%%%%%%%%%%%%%%%%%%%%%%%%%%%%%%%%%%%%%%%%

Since

\[
B_{t_{n+1}}-B_{t_n}
\sim
N(0,\Delta t),
\]

one may generate

\[
Z_n\sim N(0,1)
\]

and set

\[
\boxed{
\Delta B_n
=
\sqrt{\Delta t}\,Z_n.
}
\]

Therefore Euler--Maruyama becomes

\[
\boxed{
X_{n+1}
=
X_n
+
b(t_n,X_n)\Delta t
+
\sigma(t_n,X_n)\sqrt{\Delta t}\,Z_n.
}
\]

%%%%%%%%%%%%%%%%%%%%%%%%%%%%%%%%%%%%%%%%%%%%%%%%%%%%%%%%%%%%
\subsection{Worked Example: Euler--Maruyama}
\label{subsec:worked-euler-maruyama}
%%%%%%%%%%%%%%%%%%%%%%%%%%%%%%%%%%%%%%%%%%%%%%%%%%%%%%%%%%%%

\begin{workedexamplebox}{One Euler--Maruyama Step}

Consider

\[
dX_t
=
2X_t\,dt
+
3X_t\,dB_t.
\]

Suppose

\[
X_n=4
\]

and

\[
\Delta t=0.01.
\]

Then

\[
\Delta B_n
=
0.1Z_n.
\]

The next approximation is

\[
X_{n+1}
=
4
+
2(4)(0.01)
+
3(4)(0.1Z_n).
\]

Thus,

\begin{answerbox}
\[
\boxed{
X_{n+1}
=
4.08
+
1.2Z_n.
}
\]
\end{answerbox}

\end{workedexamplebox}

%%%%%%%%%%%%%%%%%%%%%%%%%%%%%%%%%%%%%%%%%%%%%%%%%%%%%%%%%%%%
\subsection{Euler--Maruyama for Geometric Brownian Motion}
\label{subsec:euler-maruyama-gbm}
%%%%%%%%%%%%%%%%%%%%%%%%%%%%%%%%%%%%%%%%%%%%%%%%%%%%%%%%%%%%

For

\[
dS_t
=
\mu S_t\,dt
+
\sigma S_t\,dB_t,
\]

Euler--Maruyama gives

\[
\boxed{
S_{n+1}
=
S_n
\left[
1
+
\mu\Delta t
+
\sigma\sqrt{\Delta t}\,Z_n
\right].
}
\]

This approximates the exact geometric Brownian motion solution.

For sufficiently large time steps, however, the numerical approximation
can become negative even though the true solution remains positive.

%%%%%%%%%%%%%%%%%%%%%%%%%%%%%%%%%%%%%%%%%%%%%%%%%%%%%%%%%%%%
\subsection{Milstein Method}
\label{subsec:milstein-method}
%%%%%%%%%%%%%%%%%%%%%%%%%%%%%%%%%%%%%%%%%%%%%%%%%%%%%%%%%%%%

For a scalar SDE

\[
dX_t
=
b(X_t)\,dt
+
\sigma(X_t)\,dB_t,
\]

the Milstein method improves Euler--Maruyama by including an It\^o
correction:

\[
\boxed{
\begin{aligned}
X_{n+1}
={}&
X_n
+
b(X_n)\Delta t
+
\sigma(X_n)\Delta B_n
\\
&+
\frac12
\sigma(X_n)\sigma'(X_n)
\left[
(\Delta B_n)^2-\Delta t
\right].
\end{aligned}
}
\]

The extra term reflects Brownian quadratic variation.

%%%%%%%%%%%%%%%%%%%%%%%%%%%%%%%%%%%%%%%%%%%%%%%%%%%%%%%%%%%%
\subsection{Additive Noise and Milstein}
\label{subsec:additive-noise-milstein}
%%%%%%%%%%%%%%%%%%%%%%%%%%%%%%%%%%%%%%%%%%%%%%%%%%%%%%%%%%%%

If

\[
\sigma(X_t)=\sigma
\]

is constant, then

\[
\sigma'(x)=0.
\]

Therefore the Milstein correction vanishes.

In this case,

\[
\boxed{
\text{Milstein}
=
\text{Euler--Maruyama}.
}
\]

%%%%%%%%%%%%%%%%%%%%%%%%%%%%%%%%%%%%%%%%%%%%%%%%%%%%%%%%%%%%
\subsection{Strong Numerical Convergence}
\label{subsec:strong-numerical-convergence}
%%%%%%%%%%%%%%%%%%%%%%%%%%%%%%%%%%%%%%%%%%%%%%%%%%%%%%%%%%%%

A numerical approximation \(X_T^{(\Delta)}\) converges strongly if it
approximates the actual sample path.

A typical error measure is

\[
\boxed{
\mathbb{E}
\left[
|X_T-X_T^{(\Delta)}|
\right].
}
\]

If this behaves like

\[
C(\Delta t)^p,
\]

the method has strong order \(p\).

%%%%%%%%%%%%%%%%%%%%%%%%%%%%%%%%%%%%%%%%%%%%%%%%%%%%%%%%%%%%
\subsection{Weak Numerical Convergence}
\label{subsec:weak-numerical-convergence}
%%%%%%%%%%%%%%%%%%%%%%%%%%%%%%%%%%%%%%%%%%%%%%%%%%%%%%%%%%%%

Weak convergence concerns expectations of functions of the numerical
solution.

For suitable \(f\), one studies

\[
\boxed{
\left|
\mathbb{E}[f(X_T)]
-
\mathbb{E}
[
f(X_T^{(\Delta)})
]
\right|.
}
\]

Weak convergence is important when the goal is estimating distributions
or expected payoffs rather than reproducing individual paths.

%%%%%%%%%%%%%%%%%%%%%%%%%%%%%%%%%%%%%%%%%%%%%%%%%%%%%%%%%%%%
\subsection{Euler--Maruyama Convergence Orders}
\label{subsec:euler-maruyama-orders}
%%%%%%%%%%%%%%%%%%%%%%%%%%%%%%%%%%%%%%%%%%%%%%%%%%%%%%%%%%%%

Under standard regularity conditions, Euler--Maruyama typically has:

\[
\boxed{
\text{strong order } \frac12
}
\]

and

\[
\boxed{
\text{weak order }1.
}
\]

Milstein typically improves the scalar strong order to

\[
\boxed{
1.
}
\]

%%%%%%%%%%%%%%%%%%%%%%%%%%%%%%%%%%%%%%%%%%%%%%%%%%%%%%%%%%%%
\subsection{Why SDE Numerical Accuracy Differs from ODE Accuracy}
\label{subsec:sde-vs-ode-numerics}
%%%%%%%%%%%%%%%%%%%%%%%%%%%%%%%%%%%%%%%%%%%%%%%%%%%%%%%%%%%%

For ordinary differential equations, Euler's method has deterministic
first-order behavior.

For Brownian-driven SDEs, increments scale like

\[
\sqrt{\Delta t}.
\]

This rough stochastic forcing changes numerical error accumulation.

Therefore numerical SDE methods require a separate convergence theory.

%%%%%%%%%%%%%%%%%%%%%%%%%%%%%%%%%%%%%%%%%%%%%%%%%%%%%%%%%%%%
\subsection{Stability of Numerical SDEs}
\label{subsec:sde-numerical-stability}
%%%%%%%%%%%%%%%%%%%%%%%%%%%%%%%%%%%%%%%%%%%%%%%%%%%%%%%%%%%%

A method may converge as

\[
\Delta t\to0
\]

yet behave poorly for practical finite time steps.

Stability asks whether numerical trajectories reproduce qualitative
features such as:

\[
\text{mean reversion},
\]

\[
\text{bounded moments},
\]

or

\[
\text{long-run distributions}.
\]

This is especially important for stiff SDEs.

%%%%%%%%%%%%%%%%%%%%%%%%%%%%%%%%%%%%%%%%%%%%%%%%%%%%%%%%%%%%
\subsection{Linear Stability Test SDE}
\label{subsec:linear-stability-sde}
%%%%%%%%%%%%%%%%%%%%%%%%%%%%%%%%%%%%%%%%%%%%%%%%%%%%%%%%%%%%

A common stochastic test equation is

\[
\boxed{
dX_t
=
\lambda X_t\,dt
+
\mu X_t\,dB_t.
}
\]

Its exact solution is geometric Brownian motion:

\[
\boxed{
X_t
=
X_0
\exp
\left[
\left(
\lambda-\frac12\mu^2
\right)t
+
\mu B_t
\right].
}
\]

This model is used to study mean-square and almost-sure numerical
stability.

%%%%%%%%%%%%%%%%%%%%%%%%%%%%%%%%%%%%%%%%%%%%%%%%%%%%%%%%%%%%
\subsection{Mean-Square Stability}
\label{subsec:mean-square-stability}
%%%%%%%%%%%%%%%%%%%%%%%%%%%%%%%%%%%%%%%%%%%%%%%%%%%%%%%%%%%%

A process is mean-square stable toward zero if

\[
\boxed{
\mathbb{E}[|X_t|^2]
\to0.
}
\]

For numerical methods, one asks whether the discrete approximation has
the same long-run decay property.

This is a stochastic analogue of stability analysis for ODE solvers.

%%%%%%%%%%%%%%%%%%%%%%%%%%%%%%%%%%%%%%%%%%%%%%%%%%%%%%%%%%%%
\subsection{Boundary Behavior}
\label{subsec:sde-boundary-behavior}
%%%%%%%%%%%%%%%%%%%%%%%%%%%%%%%%%%%%%%%%%%%%%%%%%%%%%%%%%%%%

The diffusion coefficient may strongly affect whether a process can
reach a boundary.

For example,

\[
\sigma(x)
\]

may vanish at \(x=0\).

Then the process may:

\[
\text{remain positive},
\]

\[
\text{hit zero},
\]

or

\[
\text{be absorbed or reflected},
\]

depending on the model.

Boundary classification is an important topic in one-dimensional
diffusion theory.

%%%%%%%%%%%%%%%%%%%%%%%%%%%%%%%%%%%%%%%%%%%%%%%%%%%%%%%%%%%%
\subsection{Absorbing Boundaries}
\label{subsec:absorbing-boundaries}
%%%%%%%%%%%%%%%%%%%%%%%%%%%%%%%%%%%%%%%%%%%%%%%%%%%%%%%%%%%%

A boundary is absorbing if once the process reaches it, the process
remains there.

For example, in a population model,

\[
X_t=0
\]

may represent extinction.

An absorbing model may satisfy

\[
\boxed{
X_t=0
\Longrightarrow
X_s=0
\quad
\text{for all }s\geq t.
}
\]

%%%%%%%%%%%%%%%%%%%%%%%%%%%%%%%%%%%%%%%%%%%%%%%%%%%%%%%%%%%%
\subsection{Reflecting Boundaries}
\label{subsec:reflecting-boundaries}
%%%%%%%%%%%%%%%%%%%%%%%%%%%%%%%%%%%%%%%%%%%%%%%%%%%%%%%%%%%%

A reflecting boundary prevents a process from leaving a permitted
region.

For instance, one may constrain a diffusion to

\[
[0,\infty).
\]

Reflected diffusions often require local-time terms such as

\[
\boxed{
dX_t
=
b(X_t)\,dt
+
\sigma(X_t)\,dB_t
+
dL_t.
}
\]

%%%%%%%%%%%%%%%%%%%%%%%%%%%%%%%%%%%%%%%%%%%%%%%%%%%%%%%%%%%%
\subsection{Stratonovich SDEs}
\label{subsec:stratonovich-sdes}
%%%%%%%%%%%%%%%%%%%%%%%%%%%%%%%%%%%%%%%%%%%%%%%%%%%%%%%%%%%%

An SDE may also be written in Stratonovich form:

\[
\boxed{
dX_t
=
a(X_t)\,dt
+
b(X_t)\circ dB_t.
}
\]

The equivalent It\^o equation is

\[
\boxed{
dX_t
=
\left[
a(X_t)
+
\frac12
b(X_t)b'(X_t)
\right]dt
+
b(X_t)\,dB_t.
}
\]

Thus the interpretation of the stochastic integral affects the drift.

%%%%%%%%%%%%%%%%%%%%%%%%%%%%%%%%%%%%%%%%%%%%%%%%%%%%%%%%%%%%
\subsection{Why Stratonovich SDEs Appear in Physics}
\label{subsec:stratonovich-physics}
%%%%%%%%%%%%%%%%%%%%%%%%%%%%%%%%%%%%%%%%%%%%%%%%%%%%%%%%%%%%

Stratonovich equations transform under smooth coordinate changes using
the ordinary chain rule.

They also arise naturally as limits of systems driven by rapidly
fluctuating smooth noise.

For these reasons, Stratonovich notation is common in physical and
geometric models.

%%%%%%%%%%%%%%%%%%%%%%%%%%%%%%%%%%%%%%%%%%%%%%%%%%%%%%%%%%%%
\subsection{Stochastic Differential Equations and Control}
\label{subsec:sde-control}
%%%%%%%%%%%%%%%%%%%%%%%%%%%%%%%%%%%%%%%%%%%%%%%%%%%%%%%%%%%%

A controlled diffusion may be written

\[
\boxed{
dX_t
=
b(t,X_t,u_t)\,dt
+
\sigma(t,X_t,u_t)\,dB_t,
}
\]

where

\[
u_t
\]

is a control process.

The goal may be to minimize an expected cost such as

\[
\boxed{
\mathbb{E}
\left[
\int_0^T
c(t,X_t,u_t)\,dt
+
g(X_T)
\right].
}
\]

This leads to stochastic optimal control.

%%%%%%%%%%%%%%%%%%%%%%%%%%%%%%%%%%%%%%%%%%%%%%%%%%%%%%%%%%%%
\subsection{Hamilton--Jacobi--Bellman Equation Preview}
\label{subsec:hjb-preview}
%%%%%%%%%%%%%%%%%%%%%%%%%%%%%%%%%%%%%%%%%%%%%%%%%%%%%%%%%%%%

The value function of a stochastic control problem often satisfies a
Hamilton--Jacobi--Bellman equation of the form

\[
\boxed{
V_t
+
\inf_u
\left\{
c
+
bV_x
+
\frac12\sigma^2V_{xx}
\right\}
=
0.
}
\]

Thus stochastic control connects optimization with second-order PDEs.

%%%%%%%%%%%%%%%%%%%%%%%%%%%%%%%%%%%%%%%%%%%%%%%%%%%%%%%%%%%%
\subsection{Risk-Neutral Dynamics Preview}
\label{subsec:risk-neutral-dynamics-preview}
%%%%%%%%%%%%%%%%%%%%%%%%%%%%%%%%%%%%%%%%%%%%%%%%%%%%%%%%%%%%

In mathematical finance, a change of probability measure can alter the
drift of an asset-price SDE.

Under a risk-neutral measure, a stylized asset may satisfy

\[
\boxed{
dS_t
=
rS_t\,dt
+
\sigma S_t\,dW_t,
}
\]

where \(r\) is the risk-free rate.

This idea relies on Girsanov's theorem and martingale pricing.

%%%%%%%%%%%%%%%%%%%%%%%%%%%%%%%%%%%%%%%%%%%%%%%%%%%%%%%%%%%%
\subsection{Black--Scholes PDE Preview}
\label{subsec:black-scholes-pde-preview}
%%%%%%%%%%%%%%%%%%%%%%%%%%%%%%%%%%%%%%%%%%%%%%%%%%%%%%%%%%%%

If an asset satisfies geometric Brownian motion under a risk-neutral
measure, It\^o calculus and hedging lead to the Black--Scholes PDE:

\[
\boxed{
V_t
+
\frac12
\sigma^2S^2V_{SS}
+
rSV_S
-
rV
=
0.
}
\]

This is one of the classic applications of SDEs.

A full treatment belongs to mathematical finance.

%%%%%%%%%%%%%%%%%%%%%%%%%%%%%%%%%%%%%%%%%%%%%%%%%%%%%%%%%%%%
\subsection{Stochastic Volatility Preview}
\label{subsec:stochastic-volatility-preview}
%%%%%%%%%%%%%%%%%%%%%%%%%%%%%%%%%%%%%%%%%%%%%%%%%%%%%%%%%%%%

Instead of assuming constant volatility, one may model volatility itself
as a stochastic process.

A generic model is

\[
\boxed{
dS_t
=
\mu S_t\,dt
+
\sqrt{V_t}S_t\,dB_t^{(1)},
}
\]

together with

\[
\boxed{
dV_t
=
a(V_t)\,dt
+
b(V_t)\,dB_t^{(2)}.
}
\]

The Brownian motions may be correlated.

This produces a multidimensional SDE system.

%%%%%%%%%%%%%%%%%%%%%%%%%%%%%%%%%%%%%%%%%%%%%%%%%%%%%%%%%%%%
\subsection{Stochastic Epidemic Models Preview}
\label{subsec:stochastic-epidemic-sdes}
%%%%%%%%%%%%%%%%%%%%%%%%%%%%%%%%%%%%%%%%%%%%%%%%%%%%%%%%%%%%

A deterministic epidemic model may be perturbed by random fluctuations.

For example,

\[
dS_t
=
-\beta S_tI_t\,dt
+
\sigma_S(S_t,I_t)\,dB_t^{(1)},
\]

\[
dI_t
=
(\beta S_tI_t-\gamma I_t)\,dt
+
\sigma_I(S_t,I_t)\,dB_t^{(2)}.
\]

Such models represent environmental or demographic uncertainty.

%%%%%%%%%%%%%%%%%%%%%%%%%%%%%%%%%%%%%%%%%%%%%%%%%%%%%%%%%%%%
\subsection{Chemical Langevin Equation Preview}
\label{subsec:chemical-langevin-preview}
%%%%%%%%%%%%%%%%%%%%%%%%%%%%%%%%%%%%%%%%%%%%%%%%%%%%%%%%%%%%

Chemical reaction systems with large but finite molecular counts can be
approximated by SDEs.

A generic chemical Langevin equation has the form

\[
\boxed{
d\mathbf{X}_t
=
\sum_{j}
\boldsymbol{\nu}_j
a_j(\mathbf{X}_t)\,dt
+
\sum_j
\boldsymbol{\nu}_j
\sqrt{
a_j(\mathbf{X}_t)
}
\,dB_t^{(j)}.
}
\]

Here:

\[
\boldsymbol{\nu}_j
\]

is the stoichiometric change vector and

\[
a_j
\]

is the reaction propensity.

%%%%%%%%%%%%%%%%%%%%%%%%%%%%%%%%%%%%%%%%%%%%%%%%%%%%%%%%%%%%
\subsection{SDE Modeling Workflow}
\label{subsec:sde-modeling-workflow}
%%%%%%%%%%%%%%%%%%%%%%%%%%%%%%%%%%%%%%%%%%%%%%%%%%%%%%%%%%%%

When constructing an SDE model:

\begin{enumerate}
    \item identify the state variable;
    \item identify deterministic dynamics;
    \item choose the source of randomness;
    \item determine whether noise should be additive or multiplicative;
    \item specify the drift \(b\);
    \item specify the diffusion coefficient \(\sigma\);
    \item determine the initial condition;
    \item check state-space constraints;
    \item determine whether boundaries should be absorbing or reflecting;
    \item verify existence and uniqueness assumptions when possible;
    \item search for an exact transformation or explicit solution;
    \item otherwise select an appropriate numerical method;
    \item validate the model against qualitative and statistical behavior.
\end{enumerate}

%%%%%%%%%%%%%%%%%%%%%%%%%%%%%%%%%%%%%%%%%%%%%%%%%%%%%%%%%%%%
\subsection{Common Errors}
\label{subsec:sde-common-errors}
%%%%%%%%%%%%%%%%%%%%%%%%%%%%%%%%%%%%%%%%%%%%%%%%%%%%%%%%%%%%

\subsubsection{Treating an SDE as an Ordinary Differential Equation}

The term

\[
dB_t
\]

cannot be manipulated as if Brownian motion were differentiable.

The SDE must be interpreted through stochastic integration.

%%%%%%%%%%%%%%%%%%%%%%%%%%%%%%%%%%%%%%%%%%%%%%%%%%%%%%%%%%%%
\subsubsection{Forgetting the It\^o Correction}

When transforming an SDE, ordinary chain rule is insufficient.

The second-order term

\[
\frac12
\sigma^2f_{xx}\,dt
\]

must be included.

%%%%%%%%%%%%%%%%%%%%%%%%%%%%%%%%%%%%%%%%%%%%%%%%%%%%%%%%%%%%
\subsubsection{Confusing Drift with Mean Path}

The drift coefficient describes local conditional mean change.

It does not generally equal the derivative of

\[
\mathbb{E}[X_t]
\]

without further calculation.

%%%%%%%%%%%%%%%%%%%%%%%%%%%%%%%%%%%%%%%%%%%%%%%%%%%%%%%%%%%%
\subsubsection{Confusing Diffusion with Variance}

The diffusion coefficient is \(\sigma\).

The instantaneous variance rate is

\[
\boxed{
\sigma^2.
}
\]

%%%%%%%%%%%%%%%%%%%%%%%%%%%%%%%%%%%%%%%%%%%%%%%%%%%%%%%%%%%%
\subsubsection{Assuming Every SDE Has a Unique Solution}

Existence and uniqueness require conditions.

Irregular coefficients may produce no solution or multiple solutions.

%%%%%%%%%%%%%%%%%%%%%%%%%%%%%%%%%%%%%%%%%%%%%%%%%%%%%%%%%%%%
\subsubsection{Confusing Strong and Weak Solutions}

Strong and weak refer to the type of probabilistic construction, not to
whether the solution is more or less accurate.

%%%%%%%%%%%%%%%%%%%%%%%%%%%%%%%%%%%%%%%%%%%%%%%%%%%%%%%%%%%%
\subsubsection{Forgetting the \(-\sigma^2/2\) Term in GBM}

For

\[
dS_t
=
\mu S_t\,dt
+
\sigma S_t\,dB_t,
\]

the exact exponent is

\[
\boxed{
\left(
\mu-\frac12\sigma^2
\right)t
+
\sigma B_t.
}
\]

%%%%%%%%%%%%%%%%%%%%%%%%%%%%%%%%%%%%%%%%%%%%%%%%%%%%%%%%%%%%
\subsubsection{Confusing OU Drift Parameters}

In

\[
dX_t
=
\theta(\mu-X_t)\,dt+\sigma\,dB_t,
\]

\(\mu\) is the long-run level and \(\theta\) is the speed of mean
reversion.

%%%%%%%%%%%%%%%%%%%%%%%%%%%%%%%%%%%%%%%%%%%%%%%%%%%%%%%%%%%%
\subsubsection{Simulating \(dB_t\) as \(N(0,1)\)}

For a numerical step of length \(\Delta t\),

\[
\boxed{
\Delta B
\sim
N(0,\Delta t),
}
\]

so one must use

\[
\boxed{
\sqrt{\Delta t}\,Z.
}
\]

%%%%%%%%%%%%%%%%%%%%%%%%%%%%%%%%%%%%%%%%%%%%%%%%%%%%%%%%%%%%
\subsubsection{Ignoring Numerical Positivity}

Euler--Maruyama may violate positivity even when the true SDE remains
positive.

This matters in population, variance, and financial models.

%%%%%%%%%%%%%%%%%%%%%%%%%%%%%%%%%%%%%%%%%%%%%%%%%%%%%%%%%%%%
\subsubsection{Confusing Strong and Weak Numerical Convergence}

Strong convergence approximates sample paths.

Weak convergence approximates distributions or expectations of
functions.

%%%%%%%%%%%%%%%%%%%%%%%%%%%%%%%%%%%%%%%%%%%%%%%%%%%%%%%%%%%%
\subsubsection{Using Too Large a Time Step}

An SDE discretization may become unstable or inaccurate if

\[
\Delta t
\]

is too large.

Convergence theory concerns the limit as

\[
\Delta t\to0.
\]

%%%%%%%%%%%%%%%%%%%%%%%%%%%%%%%%%%%%%%%%%%%%%%%%%%%%%%%%%%%%
\subsection{Applications}
\label{subsec:sde-applications}
%%%%%%%%%%%%%%%%%%%%%%%%%%%%%%%%%%%%%%%%%%%%%%%%%%%%%%%%%%%%

\begin{applicationbox}

\textbf{Finance.}
SDEs model asset prices, interest rates, volatility, credit risk, and
derivative-pricing dynamics.

\medskip

\textbf{Physics.}
Langevin equations model diffusion, friction, thermal noise, and random
forcing.

\medskip

\textbf{Biology.}
Population growth, epidemics, ecological systems, and biochemical
networks can be modeled with stochastic differential equations.

\medskip

\textbf{Engineering.}
Noisy mechanical, electrical, and control systems are naturally
represented by SDEs.

\medskip

\textbf{Stochastic Control.}
Controlled diffusions lead to Hamilton--Jacobi--Bellman equations.

\medskip

\textbf{Chemistry.}
Chemical Langevin equations approximate random reaction networks.

\medskip

\textbf{Environmental Science.}
Random forcing can represent uncertain climate, transport, or ecological
effects.

\medskip

\textbf{Numerical Mathematics.}
Euler--Maruyama, Milstein, and higher-order schemes approximate SDEs that
cannot be solved analytically.

\end{applicationbox}

%%%%%%%%%%%%%%%%%%%%%%%%%%%%%%%%%%%%%%%%%%%%%%%%%%%%%%%%%%%%
\subsection{Exercises}
\label{subsec:sde-exercises}
%%%%%%%%%%%%%%%%%%%%%%%%%%%%%%%%%%%%%%%%%%%%%%%%%%%%%%%%%%%%

\begin{exercise}[1]
Define a stochastic differential equation.
\end{exercise}

\begin{exercise}[1]
State the integral form of

\[
dX_t
=
b(t,X_t)\,dt
+
\sigma(t,X_t)\,dB_t.
\]
\end{exercise}

\begin{exercise}[1]
Identify the drift and diffusion coefficients in

\[
dX_t
=
(2-X_t)\,dt
+
3\,dB_t.
\]
\end{exercise}

\begin{exercise}[1]
Distinguish strong and weak solutions.
\end{exercise}

\begin{exercise}[1]
State standard Lipschitz and linear-growth conditions for strong
existence and uniqueness.
\end{exercise}

\begin{exercise}[1]
State the Euler--Maruyama method.
\end{exercise}

\begin{exercise}[1]
How is a Brownian increment simulated over a time step \(\Delta t\)?
\end{exercise}

\begin{exercise}[1]
Define strong numerical convergence.
\end{exercise}

\begin{exercise}[1]
Define weak numerical convergence.
\end{exercise}

\begin{exercise}[2]
Solve

\[
dX_t
=
4\,dt
+
2\,dB_t,
\qquad
X_0=3.
\]
\end{exercise}

\begin{exercise}[2]
Find the distribution of \(X_t\) in the previous exercise.
\end{exercise}

\begin{exercise}[2]
For

\[
dX_t
=
\mu\,dt+\sigma\,dB_t,
\]

find

\[
\mathbb{E}[X_t]
\]

and

\[
\operatorname{Var}(X_t).
\]
\end{exercise}

\begin{exercise}[2]
Solve geometric Brownian motion

\[
dS_t
=
\mu S_t\,dt
+
\sigma S_t\,dB_t.
\]
\end{exercise}

\begin{exercise}[2]
Show that geometric Brownian motion remains positive if

\[
S_0>0.
\]
\end{exercise}

\begin{exercise}[2]
Find the distribution of

\[
\ln S_t
\]

for geometric Brownian motion.
\end{exercise}

\begin{exercise}[2]
For

\[
dX_t
=
3(5-X_t)\,dt
+
2\,dB_t,
\]

identify the long-run mean and stationary variance.
\end{exercise}

\begin{exercise}[2]
Write one Euler--Maruyama step for

\[
dX_t
=
-X_t\,dt
+
X_t\,dB_t.
\]
\end{exercise}

\begin{exercise}[2]
Write one Milstein step for the same equation.
\end{exercise}

\begin{exercise}[3]
Use It\^o's formula to derive the exact solution of geometric Brownian
motion.
\end{exercise}

\begin{exercise}[3]
Derive

\[
\mathbb{E}[S_t]
=
S_0e^{\mu t}
\]

for geometric Brownian motion.
\end{exercise}

\begin{exercise}[3]
Derive the variance of geometric Brownian motion.
\end{exercise}

\begin{exercise}[3]
Solve the Ornstein--Uhlenbeck equation using an integrating factor.
\end{exercise}

\begin{exercise}[3]
Derive the mean of the Ornstein--Uhlenbeck process.
\end{exercise}

\begin{exercise}[3]
Derive its variance using the It\^o isometry.
\end{exercise}

\begin{exercise}[3]
Derive the stationary Ornstein--Uhlenbeck distribution.
\end{exercise}

\begin{exercise}[3]
For

\[
dX_t
=
b(X_t)\,dt
+
\sigma(X_t)\,dB_t,
\]

derive the generator using It\^o's formula.
\end{exercise}

\begin{exercise}[3]
Derive the Fokker--Planck equation formally for a one-dimensional
diffusion.
\end{exercise}

\begin{exercise}[3]
Show that the Fokker--Planck equation for Brownian motion reduces to

\[
p_t
=
\frac12p_{xx}.
\]
\end{exercise}

\begin{exercise}[3]
For a diffusion on \((a,b)\), derive the generator equation for the
probability of reaching \(b\) before \(a\).
\end{exercise}

\begin{exercise}[3]
Derive the boundary-value problem for an expected exit time.
\end{exercise}

\begin{exercise}[3]
For the stochastic logistic model, identify drift and diffusion
coefficients.
\end{exercise}

\begin{exercise}[3]
Convert

\[
dX_t
=
a(X_t)\,dt
+
b(X_t)\circ dB_t
\]

to It\^o form.
\end{exercise}

\begin{exercise}[3]
Explain why Euler--Maruyama may produce negative approximations for
geometric Brownian motion.
\end{exercise}

\begin{exercise}[3]
Explain the difference between strong order \(1/2\) and weak order \(1\).
\end{exercise}

\begin{exercise}[4]
Study the Picard-iteration proof of strong existence and uniqueness for
globally Lipschitz SDEs.
\end{exercise}

\begin{exercise}[4]
Investigate the Yamada--Watanabe theorem relating strong solutions,
weak solutions, and pathwise uniqueness.
\end{exercise}

\begin{exercise}[4]
Study one-dimensional diffusion boundary classification.
\end{exercise}

\begin{exercise}[4]
Investigate scale functions and speed measures for one-dimensional
diffusions.
\end{exercise}

\begin{exercise}[4]
Study the Feller property of diffusion semigroups.
\end{exercise}

\begin{exercise}[4]
Prove a Feynman--Kac representation for a simple diffusion.
\end{exercise}

\begin{exercise}[4]
Derive the Black--Scholes PDE from geometric Brownian motion and a
self-financing hedge.
\end{exercise}

\begin{exercise}[4]
Study the CIR process and conditions for strict positivity.
\end{exercise}

\begin{exercise}[4]
Investigate stochastic volatility models.
\end{exercise}

\begin{exercise}[4]
Study multidimensional Euler--Maruyama and correlated Brownian
increments.
\end{exercise}

\begin{exercise}[4]
Investigate higher-order stochastic numerical methods.
\end{exercise}

\begin{exercise}[4]
Study mean-square stability of Euler--Maruyama.
\end{exercise}

\begin{exercise}[4]
Investigate stochastic Runge--Kutta methods.
\end{exercise}

\begin{exercise}[4]
Study reflected SDEs and Skorokhod problems.
\end{exercise}

\begin{exercise}[4]
Investigate stochastic control and Hamilton--Jacobi--Bellman equations.
\end{exercise}

%%%%%%%%%%%%%%%%%%%%%%%%%%%%%%%%%%%%%%%%%%%%%%%%%%%%%%%%%%%%
\subsection{Conceptual Summary}
\label{subsec:sde-summary}
%%%%%%%%%%%%%%%%%%%%%%%%%%%%%%%%%%%%%%%%%%%%%%%%%%%%%%%%%%%%

A stochastic differential equation has the form

\[
\boxed{
dX_t
=
b(t,X_t)\,dt
+
\sigma(t,X_t)\,dB_t.
}
\]

Its precise meaning is

\[
\boxed{
X_t
=
X_0
+
\int_0^t
b(s,X_s)\,ds
+
\int_0^t
\sigma(s,X_s)\,dB_s.
}
\]

The drift coefficient controls local conditional mean behavior.

The diffusion coefficient controls local variance.

Arithmetic Brownian motion satisfies

\[
\boxed{
X_t
=
X_0+\mu t+\sigma B_t.
}
\]

Geometric Brownian motion satisfies

\[
\boxed{
S_t
=
S_0
\exp
\left[
\left(
\mu-\frac12\sigma^2
\right)t
+
\sigma B_t
\right].
}
\]

The Ornstein--Uhlenbeck process satisfies

\[
\boxed{
dX_t
=
\theta(\mu-X_t)\,dt
+
\sigma\,dB_t,
}
\]

with stationary law

\[
\boxed{
N
\left(
\mu,
\frac{\sigma^2}{2\theta}
\right).
}
\]

For a diffusion,

\[
\boxed{
\mathcal{L}f
=
bf'
+
\frac12\sigma^2f''.
}
\]

Its density formally satisfies

\[
\boxed{
p_t
=
-
(bp)_x
+
\frac12
(\sigma^2p)_{xx}.
}
\]

Euler--Maruyama approximates an SDE by

\[
\boxed{
X_{n+1}
=
X_n
+
b(t_n,X_n)\Delta t
+
\sigma(t_n,X_n)
\sqrt{\Delta t}\,Z_n.
}
\]

Stochastic differential equations therefore combine

\[
\boxed{
\text{differential equations}
+
\text{Brownian motion}
+
\text{It\^o calculus}
+
\text{probability}
+
\text{numerical analysis}.
}
\]

%%%%%%%%%%%%%%%%%%%%%%%%%%%%%%%%%%%%%%%%%%%%%%%%%%%%%%%%%%%%
\subsection{Section Connections}
\label{subsec:sde-connections}
%%%%%%%%%%%%%%%%%%%%%%%%%%%%%%%%%%%%%%%%%%%%%%%%%%%%%%%%%%%%

\chapterconnection{
Stochastic Differential Equations apply the stochastic integration and
It\^o calculus developed in Section~\ref{sec:stochastic-calculus} to
random dynamical systems. Stochastic Processes provide Brownian motion,
Markov processes, martingales, and long-run behavior, while
Measure-Theoretic Probability supplies the rigorous probability-space
framework. The infinitesimal generator and Fokker--Planck equations
connect SDEs with partial differential equations, while numerical SDE
methods connect directly with computational mathematics. The next
section, Random Matrix Theory, shifts from time-evolving random systems
to random linear-algebraic structures and their eigenvalue behavior.
}

%%%%%%%%%%%%%%%%%%%%%%%%%%%%%%%%%%%%%%%%%%%%%%%%%%%%%%%%%%%%
% SECTION SOURCE TRACKING
%%%%%%%%%%%%%%%%%%%%%%%%%%%%%%%%%%%%%%%%%%%%%%%%%%%%%%%%%%%%

%------------------------------------------------------------
% 7.14 Stochastic Differential Equations
%------------------------------------------------------------
%
% Source basis:
%   - Standard Itô SDE theory.
%   - Standard diffusion-process theory.
%   - Standard numerical stochastic differential-equation methods.
%
% Used for:
%   - SDE definitions;
%   - integral formulation;
%   - drift and diffusion;
%   - infinitesimal mean and variance;
%   - strong and weak solutions;
%   - pathwise uniqueness;
%   - Lipschitz existence and uniqueness conditions;
%   - arithmetic Brownian motion;
%   - geometric Brownian motion;
%   - lognormal solutions;
%   - linear SDEs;
%   - Ornstein--Uhlenbeck processes;
%   - stationary Gaussian distributions;
%   - Itô transformations of SDEs;
%   - diffusion processes;
%   - infinitesimal generators;
%   - Kolmogorov backward equations;
%   - Fokker--Planck equations;
%   - stationary densities;
%   - Feynman--Kac connections;
%   - hitting probabilities;
%   - expected hitting times;
%   - multidimensional SDEs;
%   - correlated Brownian drivers;
%   - stochastic logistic growth;
%   - Vasicek and CIR previews;
%   - Langevin equations;
%   - Euler--Maruyama method;
%   - Brownian increment simulation;
%   - Milstein method;
%   - strong and weak numerical convergence;
%   - numerical stability;
%   - boundary behavior;
%   - absorbing and reflecting boundaries;
%   - Stratonovich SDEs;
%   - stochastic control preview;
%   - Hamilton--Jacobi--Bellman equations;
%   - risk-neutral dynamics;
%   - Black--Scholes preview;
%   - stochastic volatility;
%   - stochastic epidemic models;
%   - chemical Langevin equations;
%   - applications and exercises.
%
% Notes:
%   - Random Matrix Theory is developed next in Section 7.15.
%   - Numerical SDE methods are revisited in Chapter 11.
%   - Stochastic control is developed more fully in Chapter 10.
%   - Mathematical finance applications are revisited in Chapter 13.
%
%%%%%%%%%%%%%%%%%%%%%%%%%%%%%%%%%%%%%%%%%%%%%%%%%%%%%%%%%%%%